\documentclass{article}

\usepackage{amsmath,amssymb}
\usepackage{xurl}
\usepackage[hidelinks]{hyperref}
\setlength{\emergencystretch}{2em}

% Shared commands for notes in docs/paper-gaps/.

\DeclareUnicodeCharacter{2080}{\ifmmode{}_{0}\else\textsubscript{0}\fi}
\DeclareUnicodeCharacter{2081}{\ifmmode{}_{1}\else\textsubscript{1}\fi}
\DeclareUnicodeCharacter{2082}{\ifmmode{}_{2}\else\textsubscript{2}\fi}
\DeclareUnicodeCharacter{2099}{\ifmmode{}_{n}\else\textsubscript{n}\fi}

\newcommand{\Lean}{\textsc{Lean}}
\ExplSyntaxOn
\NewDocumentCommand{\leanid}{m}{
  \ifmmode
    \textsf{\detokenize{#1}}
  \else
    \group_begin:
    \urlstyle{sf}
    \tl_set:Nn \l_tmpa_tl {#1}
    \tl_replace_all:Nnn \l_tmpa_tl {\_} {_}
    \exp_args:NV \path \l_tmpa_tl
    \group_end:
  \fi
}
\ExplSyntaxOff
\providecommand{\tr}{\operatorname{tr}}
\newcommand{\tnleanIssueBase}{https://github.com/LionSR/TNLean/issues/}
\newcommand{\tnleanPrBase}{https://github.com/LionSR/TNLean/pull/}
\newcommand{\tnleanDocBase}{https://sirui-lu.com/TNLean/docs/}
\newcommand{\ghissue}[1]{\href{\tnleanIssueBase#1}{GitHub issue \##1}}
\newcommand{\ghpr}[1]{\href{\tnleanPrBase#1}{PR \##1}}
\newcommand{\doclink}[1]{\href{\tnleanDocBase#1.html}{\path{#1}}}

% Dirac notation and the identity operator, as in blueprint/src/macros/common.tex.
% \providecommand keeps note-local definitions authoritative.
\usepackage{dsfont}
\providecommand{\ket}[1]{|#1\rangle}
\providecommand{\bra}[1]{\langle #1|}
\providecommand{\braket}[2]{\langle #1 | #2 \rangle}
\providecommand{\ketbra}[2]{|#1\rangle\!\langle #2|}
\providecommand{\Id}{\mathds{1}}
\providecommand{\one}{\mathds{1}}
\providecommand{\C}{\mathbb{C}}
\providecommand{\MN}[1]{M_{#1}(\C)}
\providecommand{\MD}{M_D(\C)}

% External referencing for paper sources.  The build helper writes sanitized
% auxiliary files under build/paper-aux/ and excludes duplicated source labels.
\usepackage{xr}

% Import a generated paper auxiliary file with a caller-supplied label prefix.
% The candidate paths support builds from docs/paper-gaps/, the repository root,
% and build/paper-aux/ itself.
\newcommand{\importpaperaux}[2]{%
  \IfFileExists{../../build/paper-aux/#2.aux}{%
    \externaldocument[#1]{../../build/paper-aux/#2}%
  }{%
    \IfFileExists{build/paper-aux/#2.aux}{%
      \externaldocument[#1]{build/paper-aux/#2}%
    }{%
      \IfFileExists{#2.aux}{%
        \externaldocument[#1]{#2}%
      }{}%
    }%
  }%
}

% General external reference macros
\newcommand{\paperref}[2]{\ref{#1-#2}}
\newcommand{\papereqref}[2]{(\ref{#1-#2})}

% CPSV16 source (1606.00608 / MPDO-22-12-17-2.tex) external document and shortcuts
\importpaperaux{cpsv16-}{1606.00608/MPDO-22-12-17-2-tnlean}
\newcommand{\cpsvref}[1]{\paperref{cpsv16}{#1}}
\newcommand{\cpsveqref}[1]{\papereqref{cpsv16}{#1}}

% Verdict marker: \gapnote{<kind>}{<status>}. Kind is the policy
% classification (clarification, local-correction, scope-restriction,
% unfaithful, false-source, open-gap); status is open, wip (elimination
% underway), resolved, or historical. Machine-read by the site index;
% prints a verdict line.
\newcommand{\gapnote}[2]{\par\noindent\textbf{Verdict:} #1 (#2).\par}


\title{Common Sector Relabeling as an Unconditional Identity}
\author{}
\date{2026-05-10}

\begin{document}
\maketitle
\gapnote{clarification}{resolved}

\section{The assertion}

Fix a physical dimension $d$ and a blocking period $p$. The physical index set
of the blocked tensor is canonically identified with the set of words of length
$p$:
\begin{align}
    \operatorname{Fin}(d^p)
        &\simeq
        (\operatorname{Fin}p\to\operatorname{Fin}d).
    \label{eq:cpsv16_sector_relabeling_hypothesis_blocked_word_equivalence}
\end{align}
The cyclic-sector construction applies this canonical multi-index relabeling
to each block.

Write $X\sim_{\mathrm{MPV}}Y$ when tensors $X$ and $Y$ generate the same
matrix product vector family at every length and for every physical word. Let
$A_k^{[p]}$ be the length-$p$ blocking of the $k$th primitive block, and let
$R(A_k^{[p]})$ denote the same blocked tensor expressed through the canonical
word relabeling in
Eq.~\ref{eq:cpsv16_sector_relabeling_hypothesis_blocked_word_equivalence}.

For scalar weights $\mu_k$, assembling the weighted block-diagonal tensor
before relabeling produces the same MPV family as assembling it afterward:
\begin{align}
    \bigoplus_k \mu_k^p\cdot A_k^{[p]}
        &\sim_{\mathrm{MPV}}
        \bigoplus_k \mu_k^p\cdot R(A_k^{[p]}).
    \label{eq:cpsv16_sector_relabeling_hypothesis_common_relabeling_identity}
\end{align}

\section{Relation to the source}

The blocked-word identification is not a separate hypothesis in
Ref.~\cite{Cirac2016MPDO_arXiv}. Blocking a tensor and then reading its physical
index through the explicit grouping of length-$p$ words are two descriptions
of the same canonical blocking operation. The source therefore uses this
identification without naming it as an additional assumption.

Equation~\ref{eq:cpsv16_sector_relabeling_hypothesis_common_relabeling_identity}
does not assert a new equivalence between independently chosen sector
coordinates. It applies one common, canonical reindexing to every block in the
family. The scalar weights and block-diagonal assembly are unchanged.

\section{Formal status}

The formalization uses the unconditional identity for the nonzero part of a
common cyclic-sector family. It does not introduce a relabeling
hypothesis.\footnote{The family-level declaration is
\leanid{CommonBlockedCyclicSectorFamily.reindexed_nonzero_part}; the common
primitive-block theorem applies it directly on both sides. The declarations
are in
\doclink{TNLean/MPS/CanonicalForm/SectorComparison/CommonBlockedCyclicSectorFamily.lean}
and
\doclink{TNLean/MPS/CanonicalForm/SectorComparison/CommonSectorTransport.lean}.}

\section{Conclusion}

The common sector relabeling is an unconditional coordinate identity.
Equation~\ref{eq:cpsv16_sector_relabeling_hypothesis_common_relabeling_identity}
formalizes the blocked-index identification used implicitly in
Ref.~\cite{Cirac2016MPDO_arXiv}; it adds no mathematical hypothesis and creates
no deviation from the source.

\bibliographystyle{alpha}
\bibliography{references}

\end{document}
